\chapter{Melléklet}

\section{A mérési konfigurációk}
\label{sec:konfiguraciok}
\paragraph{}A \ref{tab:konfiguraciok}.~táblázat minden olyan futás beállításait tartalmazza, amelynek eredménye a dolgozatban szerepel. A táblázatot a \texttt{config\_table.py} program állítja elő, API-hívás nélkül. Az adathalmazt, a modellt, a promptváltozatot, a bemenetet és az ételek számát a futás hívásonkénti rekordjaiból (\texttt{results.jsonl}) olvassa ki, a temperature-, a reasoning- és a tokenkeret-beállítást pedig a futás manifestjéből tölti be. A futás neve az eredményeit tartalmazó mappa neve. Néhány név nem a tényleges beállítást tükrözi (például az \texttt{omnifood8k\_\allowbreak{}1image\_\allowbreak{}n50\_\allowbreak{}s42} hat képpel futott), ezért a bemenet oszlop a hívásrekordok szerinti értéket mutatja. Minden képes futásban a kép hosszabbik oldala legfeljebb 1024 képpont volt, és a modell JSON-séma szerint válaszolt.

\paragraph{}A mérések során egyes paramétereket a parancssor mellett a programfájlban is módosítottam, ezért a manifestben rögzített program-ellenőrzőösszeg futásonként eltérhet a közzétett végső változatétól. A futások beállításait ettől függetlenül a manifest (a prompt szövege, a JSON-séma, a modellbeállítások, az adathalmaz-szűrés és a mintavétel) és a hívásrekordok rögzítik; ezekből készült a táblázat. A JPEG-minőséget sem a manifest, sem a hívásrekordok nem rögzítik, ezért a táblázat nem tartalmazza. A képeket minden futásban 90-es JPEG-minőséggel kódoltam, kivéve az \texttt{n5k\_\allowbreak{}lite\_\allowbreak{}macros\_\allowbreak{}q80\_\allowbreak{}n50\_\allowbreak{}s42} és a \texttt{nutrition5k\_\allowbreak{}sol\_\allowbreak{}med\_\allowbreak{}t1000\_\allowbreak{}q80\_\allowbreak{}n50\_\allowbreak{}s42} futást, amelyeket a 80-as minőséggel kódoló \texttt{jpeg80\_benchmark.py} futtatott.

\begin{landscape}
{\linespread{1}\selectfont\scriptsize\setlength{\tabcolsep}{3pt}
\begin{longtable}{>{\raggedright\arraybackslash}p{5.2cm}lll>{\raggedright\arraybackslash}p{2.6cm}l>{\raggedright\arraybackslash}p{3.2cm}r}
\caption{A dolgozatban közölt futások beállításai. Az ételek száma a futásban értékelt különböző ételek száma.}
\label{tab:konfiguraciok} \\
\hline
Futás & Dátum & Adathalmaz & Modell & Beállítás & Prompt & Bemenet & Ételek \\
 & & & & & & & száma \\
\hline
\endfirsthead
\hline
Futás & Dátum & Adathalmaz & Modell & Beállítás & Prompt & Bemenet & Ételek \\
 & & & & & & & száma \\
\hline
\endhead
\csname @@input\endcsname contents/konfiguraciok_sorok 
\end{longtable}}
\end{landscape}
